\subsection{Three kernels}
The serial code exists in three forms, each a rung of an optimisation ladder.

\begin{itemize}
\item \textbf{S0, naive two-pass.} For each query: compute all $N$ scores, find the
maximum, then accumulate. It needs only one row of scores at a time. In double
precision it is the validation reference; in single precision (S0-fp32) it is the
performance control with the same arithmetic as the other kernels.
\item \textbf{S1 basic, blocked.} Queries are processed in tiles of $B_q$ rows and
keys in tiles of $B_k$, with the online update of \cref{eq:online} between key
tiles. The loops are the obvious scalar ones: the mask is tested for every pair
and the library exponential is called for every pair.
\item \textbf{S1 tuned.} The same blocking with five changes, none of which uses a
library (\cref{tab:tuning}). This kernel is the serial baseline $T_S$ of every
speedup in this report, and the kernel inside every parallel CPU run.
\end{itemize}

\begin{table}[t]
\centering\small
\begin{tabularx}{\textwidth}{@{}l Y Y@{}}
\toprule
Change & What is done & Why it helps\\
\midrule
Packing & Each key tile is copied once per shard into a layout with one contiguous row per dimension. & The score of one query against the whole tile becomes $d$ multiply-adds of long contiguous rows over keys; the compiler vectorises this for any vector width.\\
Mask as ranges & Keys are sorted by position, so the allowed keys of a query within a tile form at most two index ranges, found by binary search. & No branch per pair; tiles that lie entirely in the future of a query tile end the loop.\\
Eight streams per pass & Scores add eight query components per pass over the tile; the value sum adds eight keys per pass over the $d_v$ components. & The running array is loaded and stored once per eight multiply-adds, so the loop is limited by arithmetic and not by memory round trips. The partial sums are written out explicitly, so no \code{-ffast-math} is needed and results do not depend on the compiler's reassociation.\\
Aligned storage & Every tensor and scratch array starts at a multiple of 4096 bytes instead of the 16 bytes that \code{malloc} guarantees. & A 512-bit load or store covers exactly one 64-byte cache line only if its address is a multiple of 64; otherwise every access straddles two lines (see below).\\
Own exponential & Range reduction $x=k\ln2+r$, a degree-6 polynomial for $e^{r}$~\cite{cephes}, and $2^k$ built from exponent bits; arguments are $\le 0$ by construction. & The loop vectorises (the library call does not); relative error about $10^{-7}$. Ranges containing arguments below $-87$ fall back to the library function so that tiny weights underflow correctly (\cref{sec:validation}).\\
\bottomrule
\end{tabularx}
\caption{What distinguishes the tuned kernel from the basic blocked one.}
\label{tab:tuning}
\end{table}

\subsection{Run time against problem size}
\begin{figure}[t]
\centering
\figureifexists{serial_scaling.pdf}
\caption{One core of a Xeon Platinum 8452Y. (a) Run time against sequence length
for the four serial codes ($d=d_v=64$, one head). (b) Time per query--key pair
against the head dimension. (c) Useful throughput of the single-precision kernels.}
\label{fig:serial}
\end{figure}

\Cref{fig:serial}a shows the expected quadratic growth. A least-squares fit of
$\log T$ against $\log N$ for $N\ge8{,}192$ gives slopes of \n{serial.slope.tuned}
for the tuned kernel, \n{serial.slope.basic} for the basic one and
\n{serial.slope.naive} for S0-fp32; doubling the sequence quadruples the time. At
$N=\n{serial.Nmax.h}$ the tuned kernel needs \n{serial.tuned.tmax}\,s for
$1.7\times10^{10}$ pairs.

In the head dimension the cost is not simply proportional (\cref{fig:serial}b). A
pair costs $2(d+d_v)$ arithmetic operations plus one exponential and some fixed
work. For the tuned kernel the multiply-adds are cheap enough for the exponential
and the bookkeeping to show at small $d$: the fitted slope in $d$ is
\n{serial.dslope.tuned}, and from $d=16$ to $d=256$ the time per pair grows by a
factor of \n{serial.dratio.tuned} instead of 16. The scalar kernels are dominated
by the multiply-adds at every $d$ and grow slightly faster than $d$
(slope \n{serial.dslope.naive} for the naive one), because at $d=256$ the keys they
stream through for every query no longer fit in the second-level cache.

\paragraph{What the tuning buys.} At $N=16{,}384$ the naive single-precision kernel
reaches \n{serial.gflops.naive}\,GFLOP/s, the basic blocked kernel
\n{serial.gflops.basic}, and the tuned kernel \n{serial.gflops.tuned}, a factor of
\n{serial.speedup.tuned} over the basic one. The two scalar kernels do not profit
from the 512-bit setting used for the tuned kernel; compiled with the compiler's
default vector width they reach \n{serial.gflops.naive.w256} and
\n{serial.gflops.basic.w256}\,GFLOP/s, so the fairest statement is a factor of
\n{serial.speedup.fair} between the best build of the basic kernel and the tuned
one. Blocking alone gains nothing here: for these sizes the naive kernel's working
set (one row of scores and the keys it streams through) is not what limits it; the
scalar inner loops are. The double precision reference is \n{serial.fp64.ratio}
times slower than its single-precision twin, which is why precision and algorithm
are compared separately.

\paragraph{Against the hardware limit.} A Sapphire Rapids core has two 512-bit
fused multiply-add units, so at the \n{hw.freq.one}\,GHz it sustains alone its peak
is $2\times16\times2\times\n{hw.freq.one}\approx\n{serial.peak}$\,GFLOP/s in single
precision. The tuned kernel reaches \n{serial.peakpct}\% of that. Three measured
effects decide how close it gets.

\emph{Alignment of the data.} \Cref{tab:align} compares three builds of the same
source that differ only in where arrays start. With the 16-byte alignment of
\code{malloc} the kernel runs at \n{align.gflops.16}\,GFLOP/s and the processor's
counters report \n{align.split.16} blocked split loads per thousand query--key
pairs: nearly every 64-byte vector crosses a cache-line boundary. Aligned to 64
bytes or to a page it runs at \n{align.gflops.4096}\,GFLOP/s, \n{align.gain.pct}\%
faster, with $\n{align.splitcount.4096}$ such events in the whole run, and the
instructions executed per cycle rise from \n{align.ipc.16} to \n{align.ipc.4096}.

\inputifexists{report/generated/table_align.tex}

\emph{Vector width.} The compiler's default for this processor is to prefer
256-bit vectors; asking for 512-bit vectors (\code{-mprefer-vector-width=512})
raises the throughput from \n{serial.gflops.w256} to \n{serial.gflops.w512}\,GFLOP/s
at the same tile sizes (\n{serial.w512.gain.pct}\%), and all results in this report
use that setting.

\emph{Tile sizes.} The packed key tile of $B_k\times d$ numbers is read once per
query row of the tile and must stay in the 48\,KiB first-level cache together with
the row of scores being written (\cref{fig:tiles}). With $d=64$ the throughput
peaks at $B_k=\n{tiles.used.bk}$ (\n{tiles.used.kib}\,KiB) and is
\n{tiles.cliff.pct}\% lower at twice that size. Longer query tiles help, from
\n{tiles.bq8.gflops}\,GFLOP/s at $B_q=8$ upwards, which is consistent with each key
tile being brought into the cache less often per pair. All other runs use
$B_q=\n{tiles.used.bq}$ and $B_k=\n{tiles.used.bk}$, which is within
\n{tiles.used.below.pct}\% of the best combination measured
($B_q=\n{tiles.best.bq}$) and leaves more query tiles per thread for the hybrid
runs.

\begin{figure}[t]
\centering
\figureifexists{serial_tiles.pdf}
\caption{Useful GFLOP/s of the tuned kernel for every combination of tile sizes
($N=16{,}384$, one core).}
\label{fig:tiles}
\end{figure}

\subsection{Two measurement pitfalls}\label{sec:pitfalls}
Both of the following were found because a result looked wrong, and both would have
distorted every scaling curve of the next section.

\paragraph{The clock needs seconds to rise.}\label{sec:ramp}
In an early version of \cref{fig:serial}c, sequences up to 4{,}096 tokens appeared
one third slower than longer ones. The cause is not the algorithm. Repeating a
\n{ramp.ms.slow}\,ms problem 400 times in a row shows a step: for the first
\n{ramp.seconds}\,s every repetition takes \n{ramp.ms.slow}\,ms, afterwards
\n{ramp.ms.fast}\,ms, a ratio of \n{ramp.ratio}, which matches the ratio of the
\n{hw.freq.one}\,GHz turbo clock to the 2.0\,GHz base clock. The processor raises its clock only after a
few seconds of sustained load. A benchmark that warms up with one short repetition
therefore measures short configurations at the base clock and long ones at the
turbo clock, and a scaling curve built from such data has a spurious bend. Every
CPU timing in this report is taken after a warm-up of at least five seconds (the
exceptions, which measure something else, are listed in \cref{tab:environment}).

\paragraph{Where the allocator puts an array.}\label{sec:alignment}
With unaligned storage the first strong-scaling runs contained efficiencies
\emph{above} one: \n{align.Emax.16} with \n{align.Emax.P.16} ranks, on one socket and
at the same clock as the serial run. Nothing in the algorithm can produce that. The
last column of \cref{tab:align} shows the cause. Whether a vector access splits
depends on the address that the allocator returns for each array, and that depends
on the sizes and the order of the allocations made before it, hence on $P$ and on
the communication variant. Identical arithmetic then runs at different speeds in
different configurations: over the no-communication control and the two ring
variants at $P\le 8$ the efficiency ranged from \n{align.Emin.16} to
\n{align.Emax.16}. With aligned storage the same twelve configurations lie between
\n{align.Emin.4096} and \n{align.Emax.4096}. A ``superlinear speedup'' of this size
would normally be explained by caches; here the counters show that it was an
accident of memory placement, and that the comparison of communication variants in
the next section would have been meaningless without the fix.
